\section{Conclusions}
\label{sec:conclusion}

The central result of this study is the large separation between the airborne
and mechanical transfer functions for the human abdominal cavity.  At the
canonical $n=2$ resonance, the transfer-function ratio is
$H_\mathrm{mech}/H_\mathrm{air}
\approx 6.5 \times 10^6\,\si{\pascal\second\squared\per\metre}$ under the SDOF
idealisation before modal-participation correction.  For the reference pair
\SI{120}{\decibel} SPL airborne loading and
$a_\mathrm{rms} = \SI{0.1}{\metre\per\second\squared}$ RMS whole-body vibration,
this system property implies a displacement ratio of approximately
$3.3 \times 10^4$; including $\Gamma_2 \approx 0.48$ reduces this to
$\sim 1.6 \times 10^4$.  Derived from first-principles shell theory and
energy-consistent reciprocity analysis, this resolves a long-standing empirical
asymmetry: occupational whole-body vibration at \SIrange{4}{8}{\hertz} causes
well-documented gastrointestinal effects, while airborne infrasound at
comparable frequencies and extreme intensities does not.
More broadly, the same analysis provides a template for comparing excitation
pathways in acoustically compact, fluid-loaded compliant structures, where
geometry and impedance determine which forcing route can access a shared mode.

The disparity has a clean physical origin.  The abdominal cavity possesses
two distinct mode families: a breathing mode ($n = 0$) near
\SI{2500}{\hertz}, governed by the fluid bulk modulus and therefore irrelevant
at infrasonic frequencies; and flexural modes ($n \geq 2$) in the
\SIrange{4}{10}{\hertz} range that match ISO~2631 abdominal transmissibility
data.  Airborne sound at these frequencies has a wavelength exceeding
\SI{85}{\metre}, placing the abdomen deep in the Rayleigh regime
($ka \approx 0.01$).  The $(ka)^n$ geometric penalty suppresses
modal-pressure content by four orders of magnitude; compounding this, the
extremely low radiation efficiency at the air--tissue interface limits the
energy available to drive the global resonance.  An energy-consistent analysis confirms an
absorption efficiency of $\sim 3 \times 10^{-14}$: the shell is, for practical
purposes, acoustically transparent.  Reaching the cellular mechanotransduction
benchmark via airborne infrasound alone would require approximately
\SI{151}{\decibel} on the energy-consistent basis
 (\SI{135}{\decibel} using the less physical pressure-based
upper-bound modal-displacement estimate)---a level at which structural damage
to the middle ear and other
severe acoustic hazards would present more immediate concerns than any
gastrointestinal effects.

Mechanical vibration, transmitted through the skeleton and abdominal wall,
couples directly and, for the reference comparison
(\SI{120}{\decibel} SPL airborne versus
$a_\mathrm{rms} = \SI{0.1}{\metre\per\second\squared}$ RMS WBV with
$\Gamma_2 = 1$), produces displacements approximately
$3.3 \times 10^4$ times larger; including the participation correction leaves
$\sim 1.6 \times 10^4$.  This disparity is not a subtle biomechanical effect
but a straightforward consequence of long-wavelength acoustics and impedance
physics.  Monte Carlo uncertainty quantification on the same energy-consistent
basis yields a reference-pair displacement ratio with median
$1.3 \times 10^4$ and 90\% credible interval
$[2.4 \times 10^3,\; 4.6 \times 10^4]$; elastic modulus accounts for 86\% of
total variance.

The underlying resonance is genuine, but the airborne-coupling mechanism
commonly invoked in brown-note claims is not supported.  What popular culture
attributes to airborne sound is more consistently explained by mechanical vibration---a distinction
that matters considerably for occupational health standards.  Airborne
infrasound cannot excite the global cavity resonance to physiologically
relevant amplitudes; whether localised compressible inclusions
(intestinal gas pockets) provide a secondary transduction pathway at extreme
intensities remains a hypothesis for future study, as such mechanisms would
operate by local compliance rather than by the global cavity resonance
established here.

Future work should verify the boundary condition assumptions through finite
element analysis of anatomically realistic geometries, investigate the role of
bowel gas pockets as localised acoustic transducers, and extend the model into
the nonlinear regime relevant to high-amplitude mechanical excitation.
